\section*{الفصل \ref{ch:b1:diagenesis-basins} --- الدياجينيز والأحواض الرسوبية}

\begin{solution}{exo:b1:diagenesis-basins:1}
حجر رملي $0.49\,e^{-2000/3704} = 0.49\times0.583 = \num{0.286}$؛ صخر طيني $0.63\,e^{-2000/1960} = 0.63\times0.360 = \num{0.227}$. تنصف المسامية حين $e^{-z/\lambda} = 1/2$، أي عند $z = \lambda\ln 2$: \qty{2567}{m} للحجر الرملي، و\qty{1359}{m} للصخر الطيني.
\end{solution}

\begin{solution}{exo:b1:diagenesis-basins:2}
$\Delta A = \Delta E + S = 30 + 50 = \qty{80}{m}$. منه \qty{60}{m} يُملأ برسوب، ف$\Delta W = 80 - 60 = \qty{20}{m}$: يتعمّق الرفّ بمقدار \qty{20}{m}.
\end{solution}

\begin{solution}{exo:b1:diagenesis-basins:3}
$h = Y(\rho_m - \rho_w)/(\rho_m - \rho_s)$: مع \qty{2300}{kg.m^{-3}}، $1.0\times2270/1000 = \qty{2.27}{km}$؛ ومع \qty{2500}{kg.m^{-3}}، $1.0\times2270/800 = \qty{2.84}{km}$. كل متر من الرسوب يحلّ محلّ متر من \omterm{def:b1:layered-earth:layers}{الوشاح} في العمود؛ فكلما كان الرسوب أكثف قلّ النقص في الكتلة مقارنةً \omterm{def:b1:layered-earth:layers}{بالوشاح}، فيلزم منه أكثر لسدّ العجز نفسه الذي يسدّه الحفر المملوء بالماء.
\end{solution}

\begin{solution}{exo:b1:diagenesis-basins:4}
(أ)~الكالسيت $100.09/2.71 = \qty{36.93}{cm^3.mol^{-1}}$؛ الدولوميت $184.40/2.84
= \qty{64.93}{cm^3.mol^{-1}}$. (ب)~يشغل مولان من الكالسيت \qty{73.86}{cm^3}، ومول واحد من الدولوميت \qty{64.93}{cm^3}: فالصلب ينكمش بمقدار $1 - 64.93/73.86 = \qty{12.1}{\%}$. (ج)~في متر مكعب، يتحوّل \qty{0.90}{m^3} من الكالسيت إلى $0.90\times(1 - 0.121) = \qty{0.791}{m^3}$ من الدولوميت؛ فتصير المسامية $1 - 0.791 = \num{0.21}$.
\end{solution}

\begin{solution}{exo:b1:diagenesis-basins:5}
(أ)~$s = 1000 - 0.63\times1960\times(1 - e^{-1000/1960}) = \qty{506.5}{m}$؛ والمسامية المتوسطة $1 - 506.5/1000 = \num{0.493}$. (ب)~مع القمّة عند \qty{2000}{m}، يعطي حلّ $s(2000, b) = 506.5$ القيمة $b = \qty{2628.7}{m}$: فالطبقة سماكتها \qty{628.7}{m}، ومسامّيتها المتوسطة $1 - 506.5/628.7 = \num{0.194}$. (ج)~الحبيبات كما هي، فالماء المفقود هو السماكة المفقودة: $1000 - 628.7 = \qty{371}{m^3}$ من الماء لكل متر مربّع.
\end{solution}

\begin{solution}{exo:b1:diagenesis-basins:6}
(أ)~$s = 300 - 0.49\times3704\times(e^{-2500/3704} - e^{-2800/3704}) =
\qty{228.10}{m}$. (ب)~من $b_0 = 300$: $s(0, 300) = 158.80$، و$1 - \varphi(300) =
0.5481$، و$b_1 = 300 + (228.10 - 158.80)/0.5481 = \qty{426.44}{m}$؛ ثم
$s(0, 426.44) = 229.07$، و$1 - \varphi = 0.5633$، و$b_2 = 426.44 - 0.96/0.5633 =
\qty{424.73}{m}$؛ والخطوة الثالثة تغيّره بأقلّ من مليمتر. وكان الحجر الرملي سماكته \qty{424.7}{m} عند السطح. (ج)~$424.7/300 = 1.42$: فقد نحّفه الدفن بعامل \num{1.42}.
\end{solution}

\begin{solution}{exo:b1:diagenesis-basins:7}
(أ)~يحتوي الجليد $52\times10^6\times916.7/1000 = \qty{4.77e7}{km^3}$ من الماء؛ موزّعةً على \qty{361.84e6}{km^2} تعطي $\qty{0.1317}{km} = \qty{132}{m}$ من مستوى البحر. (ب)~تعطي الشواطئ القديمة \qty{134}{m}، أي \qty{2}{m} أكثر. (ج)~$4.767\times
10^7/0.134 = \qty{355.7e6}{km^2}$، أي أقلّ من اليوم ب\qty{1.7}{\%}. فحين انخفض مستوى البحر برز الرفّ القاري، فكان محيط العصر الجليدي أصغر، وأدّى نقل الحجم نفسه من الماء منه إلى خفضه أكثر.
\end{solution}

\begin{solution}{exo:b1:diagenesis-basins:8}
(أ)~$(-9.1 - (-98.1))/8 = \qty{11.1}{m}$ لكل ألف سنة، قريبًا من \qty{12}{m} المستخلصة من الشواطئ القديمة. (ب)~\qty{40}{mm} في السنة تساوي \qty{40}{m} لكل ألف سنة، أي $40/12 = 3.3$ مرّات المتوسط؛ وفي خمسمئة سنة، $0.040\times500 = \qty{20}{m}$.
\end{solution}

\begin{solution}{exo:b1:diagenesis-basins:9}
(أ)~$S = 50\times2 = \qty{100}{m}$ و$\Delta E = \qty{-40}{m}$: ف$\Delta A = \qty{60}{m}$، كلّها مملوءة بالرسوب لأن العمق لم يتغيّر. (ب)~في \num{20000} سنة لا يهبط القاع إلا $50\times0.02 = \qty{1}{m}$: فينخفض \omterm{def:b1:diagenesis-basins:eustasy}{مستوى البحر النسبي} بمقدار $120 - 1 = \qty{119}{m}$. وقاع البحر، الذي كان على عمق \qty{20}{m} تحت مستوى البحر، يقف في النهاية على $119 - 20 = \qty{99}{m}$ فوقه: فيبرز الرفّ كلّه.
\end{solution}

\begin{solution}{exo:b1:diagenesis-basins:10}
(أ)~في الطبقة الرقيقة تكون المسامية شبه ثابتة عبرها: فصلبها $s \approx h(1 - \varphi(z))$ على العمق و$s \approx h_0(1 -
\varphi_0)$ عند السطح. وبمساواتهما يكون $h_0 \approx h(1 -
\varphi(z))/(1 - \varphi_0)$. (ب)~$\varphi(3000) = 0.63\,e^{-3000/1960} =
0.136$، ف$h_0 \approx 10\times0.864/0.37 = \qty{23.3}{m}$، مقابل \qty{23.1}{m} بالضبط. (ج)~حين تُعاد الطبقة إلى السطح تصير سماكتها \qty{23}{m}، ولم تعد رقيقة بالقياس إلى اضمحلال المسامية قرب السطح: فمسامّيتها تنقص من $\varphi_0$ عند قمّتها إلى أقلّ عند قاعدتها، فتتّسع حبيباتها في أقلّ قليلًا من السماكة المحسوبة ب$\varphi_0$ في كل مكان.
\end{solution}

\begin{solution}{exo:b1:diagenesis-basins:11}
(أ)~$s = 3000 - 0.63\times1960\times(1 - e^{-3000/1960}) = \qty{2032}{m}$؛ والمسامية المتوسطة $1 - 2032/3000 = \num{0.323}$، والكثافة المتوسطة $0.323\times1030 + 0.677\times2700 = \qty{2161}{kg.m^{-3}}$. (ب)~المعامل $(3300 - 1030)/(3300 - 2161) = \num{1.99}$؛ والملء يمثّل هبوطًا محمّلًا بالماء قدره $3000/1.99 = \qty{1505}{m}$. (ج)~مع كثافة الحبيبات يكون المعامل $2270/600 = \num{3.78}$، \omterm{def:b1:diagenesis-basins:subsidence}{والهبوط} \qty{793}{m} فقط. وهذا خطأ لأن ثلث الملء ماء: فالرسوب أخفّ بكثير من حبيباته، ولا تحسب في توازن الكتل إلا كثافة الرسوب المشبع.
\end{solution}

\begin{solution}{exo:b1:diagenesis-basins:12}
(أ)~يحتوي المكعب على \qty{0.90}{m^3} $= \qty{9.0e5}{cm^3}$ من الكالسيت، أي $9.0\times10^5/36.93 = \num{24370}$ مولًا. (ب)~ماغنيسيوم واحد لكل كالسيتين: \num{12180} مولًا. (ج)~يحتوي كيلوغرام من ماء البحر على $35.165\times0.0365 =
\qty{1.284}{g}$ من الماغنيسيوم، أي $1.284/24.305 = \qty{0.0528}{mol}$؛ والتفاعل يحتاج $12180/0.0528 = \qty{2.31e5}{kg}$ من ماء البحر، أي $\qty{224}{m^3}$، وهذا أكثر من مئتي مرّة حجم الصخر، حتى لو أُخذ كل أيون الماغنيسيوم. (د)~لا يتكوّن الدولوميت على نطاق واسع إلا حيث تجري أحجام كبيرة جدًّا من الماء عبر الصخر زمنًا طويلًا، أو حيث يكون الماء محلولًا ملحيًّا أغنى في الماغنيسيوم من ماء البحر: كالرفوف التي تغسلها محاليل ملحية هابطة، مثلًا، لا في رسوب يحتفظ بماء مسامه نفسه.
\end{solution}

\begin{solution}{pb:b1:diagenesis-basins:1}
\textbf{1.} الحجر الرملي: $\varphi(0) = \num{0.490}$، $\varphi(800) =
0.49\,e^{-800/3704} = \num{0.395}$. الصخر الطيني: $\varphi(800) = \num{0.419}$،
$\varphi(2600) = \num{0.167}$. الطباشير: $\varphi(2600) = \num{0.110}$،
$\varphi(3200) = \num{0.072}$.

\textbf{2.} مع $s$ من \cref{prop:b1:diagenesis-basins:solid} و$\bar\varphi
= 1 - s/h$ (السؤال~7): الحجر الرملي \num{0.441}، والصخر الطيني \num{0.274}، والطباشير
\num{0.090}.

\textbf{3.} $h\bar\varphi$: الحجر الرملي \qty{353}{m^3}، والصخر الطيني \qty{493}{m^3}، والطباشير
\qty{54}{m^3} من الماء لكل متر مربّع.

\textbf{4.} $\bar\varphi\rho_w + (1 - \bar\varphi)\rho_g$: الحجر الرملي
$0.441\times1030 + 0.559\times2650 = \qty{1936}{kg.m^{-3}}$، والصخر الطيني
\qty{2242}{kg.m^{-3}}، والطباشير \qty{2559}{kg.m^{-3}}. ويزن العمود
$1936\times800 + 2242\times1800 + 2559\times600 = \qty{7.12e6}{kg}$ لكل متر مربّع:
كثافة متوسطة \qty{2225}{kg.m^{-3}}.

\textbf{5.} الصخري $2225\times9.81\times3200 = \qty{69.9}{MPa}$؛ الهيدروستاتيكي
$1030\times9.81\times3200 = \qty{32.3}{MPa}$؛ \omterm{def:b1:stress-strain-rheology:lithostatic}{والإجهاد الفعّال}
$69.9 - 32.3 = \qty{37.5}{MPa}$.

\textbf{6.} $0.49\,e^{-z/3704} = 0.63\,e^{-z/1960}$ يعطي $z = \ln(0.63/0.49)
/(1/1960 - 1/3704) = \qty{1046}{m}$. وتحته يكون الحجر الرملي أكثر مسامية: فالصخر
الطيني يبدأ أرطب، لكنه يفقد ماءه على مسافة أقصر.

\textbf{7.} الحجر الرملي $800 - 0.49\times3704\times(1 - e^{-800/3704}) =
\qty{447.4}{m}$؛ الصخر الطيني \qty{1306.7}{m}؛ الطباشير \qty{546.0}{m}.

\textbf{8.} \qty{2300.2}{m} من الحبيبات؛ والماء يشكّل $1 - 2300.2/3200 =
\qty{28}{\%}$ من العمود.

\textbf{9.} الحجر الرملي $447.4/15 = \qty{29.8}{m}$، والصخر الطيني $1306.7/40 =
\qty{32.7}{m}$، والطباشير $546.0/20 = \qty{27.3}{m}$ من الصلب لكل مليون سنة.

\textbf{10.} كان الصخر الطيني أسرع في الإمداد، والطباشير أبطأ، لكن الثلاثة تختلف بنحو
الخُمس فقط، وهذا أقلّ بكثير مما توحي به سماكاتها اليوم.

\textbf{11.} \qty{800}{m}: قمّتها على السطح اليوم، حيث انتهى ترسيبها لتوّه، فسماكتها عند الترسيب هي سماكتها اليوم.

\textbf{12.} بحلّ $s(0, b) = 1306.7$ للصخر الطيني ب\cref{met:b1:diagenesis-basins:decompaction} نحصل على $b = \qty{2124}{m}$: كان الصخر الطيني سماكته \qty{2124}{m} عند \qty{15}{\Ma}.

\textbf{13.} $s(0, b) = 546.0$ للطباشير يعطي \qty{1071}{m} عند \qty{55}{\Ma}.

\textbf{14.} مع قمّته عند \qty{2124}{m}، يعطي $s(2124, b) = 546.0$ القيمة $b =
\qty{2748}{m}$: كانت سماكة الطباشير \qty{624}{m}.

\textbf{15.} \qty{1071}{m} عند \qty{55}{\Ma}؛ و$2124 + 624 = \qty{2748}{m}$ عند
\qty{15}{\Ma}.

\textbf{16.} من سماكات اليوم: الحجر الرملي $800/15 = \qty{53.3}{m}$،
والصخر الطيني $1800/40 = \qty{45.0}{m}$، والطباشير $600/20 = \qty{30.0}{m}$ لكل مليون
سنة. ومن السماكات عند الترسيب: \qty{53.3}{m}، و$2124/40 =
\qty{53.1}{m}$، و$1071/20 = \qty{53.6}{m}$. والمجموعة الثانية هي ذات المعنى، أي المعدّل الذي بُني به قاع البحر، وهي تُظهر معدّلًا ثابتًا تخفيه سماكات اليوم.

\textbf{17.} $1 - 1800/2124 = \qty{15}{\%}$.

\textbf{18.} عمق الماء واحد في بداية كل فترة ونهايتها، \omterm{def:b1:diagenesis-basins:accommodation}{فالاستيعاب} المتكوّن هو نموّ العمود فوق المرجع: \qty{1071}{m} من 75 إلى \qty{55}{\Ma}، و$2748 - 1071 = \qty{1677}{m}$ من 55 إلى \qty{15}{\Ma}، و$3200 - 2748 = \qty{452}{m}$ منذ \qty{15}{\Ma}.

\textbf{19.} $S = \Delta A - \Delta E$: $1677 - (-60) = \qty{1737}{m}$ من 55 إلى \qty{15}{\Ma}؛ و$452 - (-30) = \qty{482}{m}$ منذ ذلك.

\textbf{20.} $1737/40 = \qty{43}{m}$ و$482/15 = \qty{32}{m}$ لكل مليون سنة.

\textbf{21.} $60/1737 = \qty{3.5}{\%}$ و$30/482 = \qty{6.2}{\%}$. انخفض مستوى البحر في الفترتين، فأعاقت \omterm{def:b1:diagenesis-basins:eustasy}{السكونية} الملء: كان على القاعدة أن تهبط بقدر أكبر لتصنع المكان نفسه. وصنع \omterm{def:b1:diagenesis-basins:subsidence}{الهبوط} كلّ \omterm{def:b1:diagenesis-basins:accommodation}{الاستيعاب} تقريبًا.

\textbf{22.} كان الصخر الطيني والطباشير سماكتهما \qty{2748}{m} عند \qty{15}{\Ma}، وهي اليوم \qty{2400}{m}: فانضغاطهما صنع \qty{348}{m} من المكان، أي \qty{77}{\%} من \qty{452}{m} التي ملأها الحجر الرملي.

\textbf{23.} كل وحدة تنحّفت بالدفن الذي تلاها، بمقدار يعتمد على صخريّتها وعلى مدى عمق الدفن؛ ويجب قياس المكان الذي صُنع في كل فترة بالسماكات التي كانت للوحدات حينئذٍ، ولا يعطيها إلا \omterm{def:b1:diagenesis-basins:compaction}{فكّ الانضغاط}. وحتى عندئذٍ يتضمّن \omterm{def:b1:diagenesis-basins:subsidence}{الهبوط} الذي يُوجد هكذا الغوص الناجم عن وزن الرسوب، ويجب أن يُزال للوصول إلى \omterm{def:b1:diagenesis-basins:subsidence}{الهبوط التكتوني}.

\textbf{24.} \qty{2124}{m}، نحو \qty{2.1}{km}، أي \qty{324}{m} أكثر من \qty{1800}{m} التي تجدها البئر اليوم.
\end{solution}
